% --- Tabel Deskripsi UC-02 ---
\begin{table}[H]
    \centering
    \caption{Deskripsi Use Case Mengerjakan Latihan Adaptif (UC-02)}
    \label{tbl:desc_uc02}
    \begin{tabular}{|p{0.25\textwidth}|p{0.70\textwidth}|}
        \hline
        \textbf{Nama Use Case} & \textbf{Mengerjakan Latihan Adaptif} \\
        \hline
        \textbf{Aktor} & Siswa \\
        \hline
        \textbf{FR Terkait} & FR-02, FR-08 \\
        \hline
        \textbf{Deskripsi} & Siswa mengerjakan soal latihan yang dipilihkan secara dinamis oleh sistem berdasarkan estimasi kemampuan saat ini untuk meningkatkan kompetensi. \\
        \hline
        \textbf{Kondisi sebelum} & Siswa telah login; Nilai parameter Elo rating ($\theta$) siswa tersedia. \\
        \hline
        \textbf{Kondisi sesudah} & Skor Elo siswa diperbarui; Sistem mengevaluasi tren variabilitas skor (deteksi stagnasi). \\
        \hline
        \textbf{Main Flow} & \begin{enumerate}[noitemsep,topsep=0pt,leftmargin=*]
            \item Sistem memilih latihan soal dengan tingkat kesulitan setara ($D_j \approx R_i$).
            \item Sistem menampilkan soal kepada siswa.
            \item Siswa menginputkan jawaban.
            \item Sistem memvalidasi jawaban dan menghitung skor aktual.
            \item Sistem memperbarui profil kompetensi siswa.
        \end{enumerate} \\
        \hline
    \end{tabular}
\end{table}